% Figura do referencial: bagging x boosting (desenho próprio).
\begin{figure}[H]
    \centering
    \begin{tikzpicture}[
        font=\footnotesize,
        caixa/.style={draw, rounded corners=2pt, align=center, fill=gray!8, minimum height=0.8cm, inner sep=3pt, text width=1.8cm, execute at begin node={\hyphenpenalty=10000\exhyphenpenalty=10000}},
        dado/.style={caixa, fill=white},
        saida/.style={caixa, fill=gray!25, font=\footnotesize\bfseries, text width=3.2cm},
        seta/.style={-{Stealth[length=2mm]}, thick},
        etapa/.style={font=\scriptsize\bfseries\scshape, text=gray!70!black}
    ]
        % Bagging (esquerda)
        \node[etapa] at (2.25, 1.0) {\textit{Bagging} (paralelo)};
        \node[dado, text width=2.4cm] (bbase) at (2.25, 0) {Base de treino};
        \foreach \i/\x in {1/0.0, 2/2.25, 3/4.5} {
            \node[dado] (ba\i) at (\x, -1.5) {Amostra \i};
            \node[caixa] (bm\i) at (\x, -3.0) {Árvore \i};
            \draw[seta] (bbase.south) -- (ba\i.north);
            \draw[seta] (ba\i) -- (bm\i);
        }
        \node[saida] (bout) at (2.25, -4.6) {Média ou voto};
        \foreach \i in {1,2,3} \draw[seta] (bm\i.south) -- (bout.north);

        % Separador
        \draw[gray!50, dashed] (5.85, 1.2) -- (5.85, -5.2);

        % Boosting (direita)
        \node[etapa] at (10.0, 1.0) {\textit{Boosting} (sequencial)};
        \node[dado, text width=2.4cm] (gbase) at (10.0, 0) {Base de treino};
        \node[caixa, text width=1.55cm] (gm1) at (7.3, -1.5) {Árvore 1};
        \node[caixa, text width=1.55cm] (gm2) at (10.0, -1.5) {Árvore 2};
        \node[caixa, text width=1.55cm] (gm3) at (12.7, -1.5) {Árvore 3};
        \draw[seta] (gbase.south) -- (gm1.north);
        \draw[seta] (gm1) -- node[above, font=\scriptsize] {erros} (gm2);
        \draw[seta] (gm2) -- node[above, font=\scriptsize] {erros} (gm3);
        \node[saida] (gout) at (10.0, -4.6) {Soma das contribuições};
        \foreach \i in {1,2,3} \draw[seta] (gm\i.south) -- (gout.north);
    \end{tikzpicture}
    \caption[Diferença estrutural entre \textit{Bagging} e \textit{Boosting}]{Diferença estrutural entre as duas estratégias de \textit{ensemble}. À esquerda, no \textit{Bagging}, a base de treino é reamostrada com reposição, cada amostra treina uma árvore de forma paralela e independente, e as predições são combinadas por média ou voto. À direita, no \textit{Boosting}, as árvores são treinadas em sequência: cada uma é ajustada aos erros que o conjunto das anteriores ainda comete, e as predições são somadas. O XGBoost, usado no Rating da SEFAZ-DF, pertence à segunda família. Elaborada pelos autores com base em \citeonline{hastieElementsStatisticalLearning2009}.}
    \label{fig:ensemble_comparison}
\end{figure}
